\section{Discussion}\label{sec:discussion}

\textbf{Recording with the BrainLink headset.} Each member recorded 30 blocks of 20~s (ten rounds of Relax, Focus and Stroop). Pairing the headset and selecting the serial port caused no problems. However, in one session we only noticed after all rounds had been recorded that the headset had not actually been connected, so the whole session had to be repeated. A few files also turned out to be slightly shorter than 20~s, so we recorded those rounds again. Since the first session had likely been disturbed by the crowded, noisy lab and a tired subject, we compared the spectra with the reference examples before keeping a session. The dry electrode at Fp1 is sensitive to contact quality and movement, and both the Focus and the Stroop tasks are answered with key presses. In our data, the Stroop recordings show the largest amplitude, and also the highest 90th percentile and maximum of the 1~s window amplitude, which is partly useful for classification but cannot be separated from muscle and eye artifacts with a single channel. In the first rounds of one member's session, we observed particularly large artifacts (the standard deviation of five early recordings of B12901069 was 1.7 to 2.8 times the median over all files), so the early rounds of a session seem to be less reliable. We kept all recordings, used the filtering described in Section~\ref{sec:method} and relied on the distribution statistics of the window features (mean, standard deviation, 90th percentile and maximum).

\textbf{Software challenges.} The first version of our pipeline used one Welch spectrum per segment with 21 features (relative and log band power, ratios and time-domain statistics). It reached only about 60\% LOSO accuracy at best, because the absolute and relative band powers differ strongly between members, so a model trained on two members does not transfer to the third. We addressed this with two changes: (i) the distribution of 1~s window features inside each segment, which keeps short bursts that a single spectrum averages away, and (ii) a subject-wise z-score of the features. The framework calls \texttt{extract\_features} on segments without a subject ID, so the normalization had to be implemented inside Part~2 by building a segment-to-feature table once and normalizing each subject with its own statistics. Features that looked helpful in isolation, such as relative band power alone or additional spectral ratios, did not help across subjects until the normalization was added.

\textbf{Limitations and future work.} A single dry electrode at Fp1 captures frontal activity only and mixes EEG with eye-blink and muscle signals. Alpha reactivity, which is the main cue for Relax, is much weaker at Fp1 than over the occipital cortex, and B12901057 shows only a weak alpha difference between Relax and Focus. With three subjects the results are also sensitive to a single member. Future work could (i) record more subjects and more electrode positions, (ii) control or label the eye and key-press artifacts (for example with a blink-removal step or a hands-free response), (iii) test normalization that only needs a short resting baseline from the new user, or domain-adaptation methods, to reduce the need for a recording of the target user, and (iv) evaluate on subjects not used for any design decision.
